\section{问题重述}

\subsection{研究背景与实际意义}

自然灾害尤其是突发性暴雨引发的洪涝与地质灾害，常造成灾区地面交通大面积瘫痪、供电中断及通信基站严重受损。在复杂的山地丘陵环境中，泥石流与滑坡使得救援人员和地面车辆难以在“黄金 72 小时”内抵达受灾孤岛，急救药品、血液、净水及高能量即食口粮等关键救灾物资面临运输受阻的严重困境。与此同时，山区复杂起伏的地理高程对无线电电磁波传播形成强烈的视距遮挡（Line-of-Sight, LOS），导致地面调度指挥中心与前线作业单元之间的通信链路频繁中断，极大增加了次生灾害防范难度与协同作业风险。

面对极端复杂的山地环境与严苛的应急救援时效要求，利用无人机（Unmanned Aerial Vehicles, UAVs）开展立体救援具有响应迅速、地形适应性强、机动灵活等不可替代的优势。然而，在复杂山地多旋翼无人机应急救援体系中，决策者面临着多项紧密耦合的物理与管理硬约束挑战：
\begin{enumerate}
    \item \textbf{重力势能与能量边界耦合}：多旋翼无人机爬升克服重力势能做功巨大，且其动力学能耗与机身载重呈高度非线性的三点二次方关系，使得不同起伏航段的有效航程极不对称；
    \item \textbf{机电分离与车电异步流水线管理}：机载大容量锂电池充电耗时显著长于飞行时间，需建立实体机与共享电池的高效周转排班体系，避免地面排队闲置；
    \item \textbf{地形起伏诱发的三维视距深衰落}：山区山脊高程远超视线高度，单一地面网关难以维系全空域连续可通视通信，必须统筹引入长续航空中中继节点进行空间协同覆盖；
    \item \textbf{应急责任分区隔离与资源刚性约束}：在责任分区独立封闭运行模式下，跨区装备调配被严格禁止，引发时分复用阻断与瞬时并发需求叠加。
\end{enumerate}

因此，以广西横州镇龙乡 30 米数字高程模型（DEM）与真实受灾点分布为地理基础，建立无人机运输与通信协同优化模型体系，为山区应急物资运送与通信保障提供可落地的科学决策支持。

\subsection{赛题任务拆解与研究内容}

根据赛题要求，本文在统一的 30 米 DEM 高程场、离散物资清单与异构装备参数体系下，系统拆解为四个层层递进的核心研究任务：

\textbf{任务一：单点往返运输能力评估与货箱组批优化}
针对由临时调度中心（$O01$）至 15 个受灾服务区（$S001 \sim S015$）的单点往返直达任务，要求在保证无人机返航安全电量余量（$\eta \ge 0.20$）的前提下，精确评估 A、B、C 三种异构机型的最大安全物理有效载荷；进而，在 80 件不可拆分货箱的质量与容积双重硬约束下，构建单点组批优化模型，求解最少架次与最低能耗的运输方案，并对安全电量余量 $\eta$ 进行全工况敏感性分析。

\textbf{任务二：异构无人机多点多架次协同运输调度}
放开单点直达限制，允许无人机在单架次中相继访问多个临近服务区实施回路投送。在 8 架实体机（A 型 4 架、B 型 2 架、C 型 2 架）与 14 组共享电池（A 型 6 组、B 型 4 组、C 型 4 组）资源硬约束下，综合考虑 30 件首批急救保障物资的硬时限交付要求、80 件货箱的期望时限、地面电池两阶段非线性充放电及机坪周转时序，建立多点多架次协同网络优化模型，并设计高效全局求解算法，输出精确至秒级的作业时序。

\textbf{任务三：通信约束下的运输与中继联合协同调度}
在任务二所确定的 24 个运输航次时空航迹基础上，利用 30 米 DEM 开展全航程三维视距遮挡判定与双向链路预算分析，精准定位直连通信中断时段与空域坐标；结合山体地形特征实施战略制高点空中中继选址；在 2 架实体中继机与 6 组共享能源组件的物理配置下，设计中继机巡航、悬停保活、地面换电与多机轮转调度方案，确保全航程实现 100\% 连续通信，并优化系统联合完工时间与能耗。

\textbf{任务四：应急救援任务分区与独立资源配置优化}
为降低大规模系统联合调度的组织复杂度，需将 15 个受灾服务区划分为 2 个任务组（$K=2$）和 3 个任务组（$K=3$）。基于“同一航次涉及多个服务区必须划入同一任务组”的刚性拓扑连通规则，证明分区的代数拓扑结构；在“任务执行期间装备不得跨组调配”的强隔离要求下，核算各任务组独立执行所需的运输机、共享电池、中继机及组件的峰值并发规模，量化现有库存缺口并深度剖析其管理学诱因。
